\section{Zero-Inflated GLMs for Insurance Claims}
\areastatus{Explore} \hfill \textit{Opened: \todo{date}}

\subsection*{Core Question}
How can zero-inflated and hurdle GLM extensions be adapted to model
insurance claim frequency in a market with a very high proportion of
zero-claim policyholders, and what does this reveal that a standard
Poisson/negative-binomial model misses?

\subsection*{Why This Area}
Sits at the intersection of two priority areas from the landscape scan
(Missing Data + Robust Statistics, and Statistical Machine Learning), and
insurance claims data is a realistic, locally-relevant application.

\subsection*{Key Papers So Far}
\begin{itemize}
\item Lambert (1992) -- the original zero-inflated Poisson (ZIP) model;
general statistical foundation, not insurance-specific.
\item Yip \& Yau (2005) -- first direct application of ZIP to claim
frequency; establishes the problem is real in claims data specifically.
\item Boucher, Denuit \& Guillen (2007) -- compares zero-inflated mixed
Poisson against hurdle models for claim-count risk classification.
\item Alomair (2024) -- direct comparison of ZIP/ZINB/hurdle against
ML methods (SVM, RF, ANN) on claim frequency; useful methodology
template.
\item Gu (2024) -- current frontier of ZIP-Tweedie modelling
(dispersion + gradient boosting); benchmark for existing sophistication.
\item Workie \& Azene (2021) -- adjacent-field evidence: zero-inflated/
hurdle modelling is mature in African demographic data, but not yet
applied to insurance claims in that setting. See note in
\texttt{02-literature/notes/}.
\end{itemize}

\subsection*{Literature Landscape Check (23 Aug 2026)}
Searched via FastTrack across OpenAlex/Semantic Scholar, 2020--2025
window, query ``zero-inflated GLM insurance claim frequency modeling'':
\textbf{{\raise0pt\hbox{\textasciitilde}}216 papers} in this window --
an active field, not a closing one. Two dominant live threads:
(a) ZIP/hurdle vs.\ gradient-boosting/ML comparisons, and (b)
copula/dependence-structure extensions. Venue spread across the top 20
most relevant hits: actuarial-specialist journals (\textit{Risks},
ASTIN Bulletin, NAAJ, Journal of Risk \& Insurance) accounted for 8;
general statistics/methodology journals for 6; ML/CS venues
(incl.\ 2 arXiv preprints) for 3; PLoS ONE for 2.

\textbf{Gap check result:} searching specifically for zero-inflated
models in developing-market/African insurance contexts returned almost
no insurance-claims papers -- instead surfacing a mature adjacent
literature applying the same methodology to demographic/public-health
counts (child mortality, fertility, antenatal care). This is read as
evidence the gap is genuinely open: the methodology is proven to
transfer to African/developing-market data structurally, but nobody in
this window has pointed it at insurance claims specifically.

\subsection*{Next Steps}
\begin{itemize}[label=$\square$]
\item Contact a local insurer, or identify a public claims dataset
(e.g.\ French MTPL via \texttt{CASdatasets} in R) as a fallback if
local data access doesn't come through.
\item Read Alomair (2024) and Gu (2024) in full to confirm neither
already tests small/thin-exposure sensitivity under a different framing
than the search caught.
\item Draft a one-paragraph gap statement combining Workie \& Azene
(2021) with the actuarial ZIP/hurdle literature -- this is the
strongest available novelty argument right now.
\item Re-run this same landscape check on garch-frontier-markets and
survival-censoring before committing further time here.
\end{itemize}